\SourcePage{278}{281}
\section{The indistinguishability principle for identical particles}

In classical mechanics, identical particles (electrons, for example) retain their ``individuality'' despite having identical physical properties. One may imagine assigning numbers, at some instant, to the particles comprising a given physical system and thereafter following the motion of each along its own trajectory; the particles could then be identified at any later instant.

Quantum mechanics presents an entirely different situation. We have already noted repeatedly that the uncertainty principle deprives the electron trajectory of all meaning. If an electron's position is known exactly at the present instant, its coordinates have no definite values at all an instant later. Hence localizing and numbering the electrons at one instant does nothing toward identifying them at subsequent instants. If, at another instant, an electron is localized at some point in space, it is impossible to say which particular electron has arrived there. Thus quantum mechanics rules out in principle any means of following identical particles separately and thereby telling them apart. In this sense, identical particles lose their ``individuality'' completely in quantum mechanics. Their equality in physical properties acquires a far deeper significance here: it entails their complete indistinguishability.

This \emph{indistinguishability principle} for identical particles, as it is called, is fundamental to the quantum theory of systems composed of such particles. Begin with a system containing only two particles. Since the particles are identical, two states related merely by interchanging them must be physically equivalent in every respect. Consequently, such an interchange can alter the system's wave function only by

\SourcePage{279}{282}
an immaterial phase factor. Let $\psi(\xi_1,\xi_2)$ be the system's wave function, with $\xi_1$ and $\xi_2$ denoting, conventionally, the three coordinates together with the spin projection of each particle. We must then have
\[
  \psi(\xi_1,\xi_2)=e^{i\alpha}\psi(\xi_2,\xi_1),
\]
where $\alpha$ is a real constant. Repeating the interchange restores the initial state, while multiplying $\psi$ by $e^{2i\alpha}$. It follows that $e^{2i\alpha}=1$, or $e^{i\alpha}=\pm1$. Thus $\psi(\xi_1,\xi_2)=\pm\psi(\xi_2,\xi_1)$.

There are therefore only two alternatives: the wave function is symmetric (it is wholly unchanged by interchanging the particles), or it is antisymmetric (its sign is reversed by the interchange). Plainly, the wave functions of all states of one and the same system must possess the same symmetry. Otherwise, a state formed by superposing states of different symmetries would have a wave function that was neither symmetric nor antisymmetric.

The result extends immediately to systems containing any number of identical particles. Indeed, particle identity makes it clear that, if one pair is described, say, by symmetric wave functions, every other pair of the same kind must share that property. The wave function of identical particles must accordingly either remain unchanged when any pair is interchanged (and hence under any permutation whatever), or reverse sign whenever a pair is interchanged. The former is called a \emph{symmetric} wave function and the latter an \emph{antisymmetric} wave function. Which of these two kinds of wave function describes a particle depends on its species. Particles described by antisymmetric functions are said to obey \emph{Fermi--Dirac statistics}, or to be \emph{fermions}; particles described by symmetric functions are said to obey \emph{Bose--Einstein statistics}, or to be \emph{bosons}\footnote{The terminology comes from the names of the statistics describing ideal gases made up, respectively, of particles with antisymmetric or symmetric wave functions. In fact, the distinction here is not merely between two statistics, but essentially between two mechanics. Fermi statistics was proposed by E.~Fermi for electrons in 1926, and its relation to quantum mechanics was established by Dirac (1926). Bose statistics was proposed by S.~Bose for light quanta and generalized by Einstein (1924).}.

\SourcePage{280}{283}
Relativistic quantum mechanics permits one to show (see IV, §~25) that a particle's statistics is uniquely tied to its spin: particles of half-integral spin are fermions, whereas particles of integral spin are bosons.

For composite particles the statistics is fixed by the parity of the number of elementary fermions they contain. Indeed, interchanging two identical composite particles amounts to interchanging several pairs of identical elementary particles simultaneously. Interchanging bosons leaves the wave function unchanged, while interchanging fermions reverses its sign. Composite particles containing an odd number of elementary fermions therefore obey Fermi statistics, and those containing an even number obey Bose statistics. This conclusion naturally agrees with the general rule just stated: a composite particle has integral or half-integral spin according as the number of its half-integral-spin constituents is even or odd.

Thus atomic nuclei of odd atomic weight (that is, containing an odd total number of protons and neutrons) obey Fermi statistics, while nuclei of even weight obey Bose statistics. For atoms, which contain electrons in addition to the nucleus, the statistics is evidently determined by whether the sum of the atomic weight and atomic number is even or odd.

Consider a system of $N$ identical particles whose mutual interaction may be neglected. Let $\psi_1,\psi_2,\ldots$ be wave functions for the various stationary states available to each particle separately. The state of the system as a whole can be specified by listing the state numbers occupied by its individual particles. We must determine how the wave function $\psi$ of the entire system is to be constructed from the functions $\psi_1,\psi_2,\ldots$.

Let $p_1,p_2,\ldots,p_N$ denote the state numbers of the individual particles; repetitions among these numbers are allowed. For a boson system, the wave function $\psi(\xi_1,\xi_2,\ldots,\xi_N)$ is a sum of products of the form
\begin{equation}
  \psi_{p_1}(\xi_1)\psi_{p_2}(\xi_2)\cdots\psi_{p_N}(\xi_N)
  \tag{61.1}
\end{equation}
over all distinct permutations of the indices $p_1,p_2,\ldots$. This sum plainly has the requisite symmetry. For example, for two particles occupying different states ($p_1\ne p_2$),
\begin{equation}
 \psi(\xi_1,\xi_2)=\frac1{\sqrt2}\left[\psi_{p_1}(\xi_1)\psi_{p_2}(\xi_2)+\psi_{p_1}(\xi_2)\psi_{p_2}(\xi_1)\right].
 \tag{61.2}
\end{equation}

\SourcePage{281}{284}
The factor $1/\sqrt2$ provides normalization (all the functions $\psi_1,\psi_2,\ldots$ are mutually orthogonal and are assumed normalized). More generally, for a system of any number $N$ of particles, the normalized wave function is
\begin{equation}
 \psi_{N_1N_2\ldots}=\left(\frac{N_1!N_2!\cdots}{N!}\right)^{1/2}
 \sum \psi_{p_1}(\xi_1)\psi_{p_2}(\xi_2)\cdots\psi_{p_N}(\xi_N),
 \tag{61.3}
\end{equation}
where the sum runs through all distinct permutations of $p_1,p_2,\ldots,p_N$, and $N_i$ gives the number of these indices having the value $i$ (so that $\sum N_i=N$). On integrating $|\psi|^2$ over $d\xi_1d\xi_2\cdots d\xi_N$\footnote{Here and in §§~64, 65, integration over $d\xi$ conventionally means integration over the coordinates together with summation over $\sigma$.}, every term vanishes except the squared modulus of each term in the sum. Since the number of terms in the sum (61.3) is evidently $N!/(N_1!N_2!\cdots)$, the normalization coefficient displayed there follows.

For fermions, $\psi$ is the antisymmetric combination of the products (61.1). In a two-particle system, for example,
\begin{equation}
 \psi(\xi_1,\xi_2)=\frac1{\sqrt2}\left[\psi_{p_1}(\xi_1)\psi_{p_2}(\xi_2)-\psi_{p_1}(\xi_2)\psi_{p_2}(\xi_1)\right].
 \tag{61.4}
\end{equation}
For $N$ particles in general, the system's wave function takes the determinant form
\begin{equation}
 \psi_{N_1N_2\ldots}=\frac1{\sqrt{N!}}
 \begin{vmatrix}
  \psi_{p_1}(\xi_1)&\psi_{p_1}(\xi_2)&\cdots&\psi_{p_1}(\xi_N)\\
  \psi_{p_2}(\xi_1)&\psi_{p_2}(\xi_2)&\cdots&\psi_{p_2}(\xi_N)\\
  \cdots&\cdots&\cdots&\cdots\\
  \psi_{p_N}(\xi_1)&\psi_{p_N}(\xi_2)&\cdots&\psi_{p_N}(\xi_N)
 \end{vmatrix}.
 \tag{61.5}
\end{equation}
Interchanging two particles interchanges two columns of this determinant and therefore reverses its sign.

Formula (61.5) has an important consequence. If two of the numbers $p_1,p_2,\ldots$ coincide, two rows of the determinant coincide and the determinant vanishes identically. It can be nonzero only when all the numbers $p_1,p_2,\ldots$ are different. Thus no two (or more) particles in a system of identical fermions can occupy the same state simultaneously. This is the so-called \emph{Pauli principle} (W.~Pauli, 1925).
